\documentclass[10pt]{article}

\usepackage[T1]{fontenc}
\usepackage{lmodern}
\usepackage{amsmath,amssymb,mathtools}
\usepackage[margin=0.75in]{geometry}
\usepackage{microtype}
\usepackage[hidelinks]{hyperref}

\allowdisplaybreaks[4]
\setlength{\parindent}{0pt}
\setlength{\parskip}{0.5em}
\raggedbottom

\newcommand{\Eref}{E_{\mathrm{ref}}}
\newcommand{\vac}{\Phi_0}
\newcommand{\tC}{\widetilde{C}}
\newcommand{\tPhi}{\widetilde{\Phi}}
\newcommand{\tPsi}{\widetilde{\Psi}}
\newcommand{\tsigma}{\widetilde{\sigma}}

\title{Closed-Shell Configuration-Interaction Expressions\\
       in the HJO Biorthogonal Framework}
\author{}
\date{25 November 2026}

\begin{document}

\maketitle
\begingroup
\small
\tableofcontents
\endgroup

\section{Definitions and conventions}

\subsection{Indices and summation convention}

Indices $i,j,k,l,\ldots$ denote spatial orbitals that are doubly occupied in
the closed-shell reference determinant $\lvert\vac\rangle$;
$a,b,c,d,\ldots$ denote virtual spatial orbitals; and
$p,q,r,s,\ldots$ denote general spatial orbitals. Repeated internal indices
are summed. The orbitals are orthonormal, but they are not assumed to be
canonical.

Spin labels are denoted by $\sigma,\tau\in\{\alpha,\beta\}$. Creation and
annihilation operators obey the usual fermionic anticommutation relations.
Normal ordering with respect to $\lvert\vac\rangle$ is denoted by braces.

\subsection{Spin-free excitation operators}

The unitary-group generator is
\begin{equation}
 E_{pq}
 =\sum_{\sigma}a_{p\sigma}^{\dagger}a_{q\sigma}.
\end{equation}
It satisfies
\begin{equation}
 [E_{pq},E_{rs}]
 =\delta_{qr}E_{ps}-\delta_{ps}E_{rq}.
\end{equation}
The spin-free two-electron excitation operator is
\begin{align}
 e_{pqrs}
 &\equiv E_{pr}E_{qs}-\delta_{rq}E_{ps}
 \\
 &=\sum_{\sigma\tau}
   a_{p\sigma}^{\dagger}a_{q\tau}^{\dagger}
   a_{s\tau}a_{r\sigma}.
\end{align}
Consequently,
\begin{equation}
 e_{pqrs}=e_{qpsr}.
\end{equation}

\subsection{One- and two-electron integrals}

The one-electron matrix elements are
\begin{equation}
 h^{p}_{q}
 =\int\!\mathrm{d}\mathbf r\,
   \phi_p^{*}(\mathbf r)\,\hat h\,\phi_q(\mathbf r).
\end{equation}
The two-electron integrals use physicists' ordering,
\begin{equation}
 u^{pq}_{rs}
 =\iint\!\mathrm{d}\mathbf r_1\,\mathrm{d}\mathbf r_2\,
   \phi_p^{*}(\mathbf r_1)\phi_q^{*}(\mathbf r_2)
   \hat w(1,2)
   \phi_r(\mathbf r_1)\phi_s(\mathbf r_2).
\end{equation}
For a symmetric two-particle interaction,
\begin{equation}
 u^{pq}_{rs}=u^{qp}_{sr},
 \qquad
 u^{pq}_{rs}=\bigl(u^{rs}_{pq}\bigr)^{*}.
\end{equation}
The spatial-orbital integrals $u^{pq}_{rs}$ are \emph{not}
antisymmetrized. Coulomb and exchange contributions therefore occur
explicitly. The sigma-vector expressions below use the abbreviation
\begin{equation}
 L^{pq}_{rs}\equiv 2u^{pq}_{rs}-u^{pq}_{sr}.
\end{equation}

The closed-shell reference Fock matrix and reference energy are
\begin{align}
 f^{p}_{q}
 &=h^{p}_{q}
   +\sum_i\left(2u^{pi}_{qi}-u^{pi}_{iq}\right),
 \\
 \Eref
 &=2\sum_i h^{i}_{i}
   +2\sum_{ij}u^{ij}_{ij}
   -\sum_{ij}u^{ij}_{ji}.
\end{align}

\subsection{Hamiltonian}

The spin-free electronic Hamiltonian is
\begin{equation}
 H
 =\sum_{pq}h^{p}_{q}E_{pq}
  +\frac12\sum_{pqrs}u^{pq}_{rs}e_{pqrs}.
\end{equation}
Normal ordering with respect to the closed-shell reference gives
\begin{equation}
 H
 =\Eref
  +\sum_{pq}f^{p}_{q}\{E_{pq}\}
  +\frac12\sum_{pqrs}u^{pq}_{rs}\{e_{pqrs}\}.
\end{equation}

\subsection{HJO excitation and projection operators}

The right excitation operators used for singles and redundant doubles are
\begin{align}
 X_i^a&=E_{ai},
 &
 X_{ij}^{ab}&=E_{ai}E_{bj}.
\end{align}
The HJO biorthogonal projection operators are
\begin{align}
 Y_a^i
 &=\frac12E_{ia},
 \\
 Y_{ab}^{ij}
 &=\frac13E_{jb}E_{ia}+\frac16E_{ib}E_{ja}.
\end{align}
They define
\begin{align}
 \lvert\Phi_i^a\rangle
 &=X_i^a\lvert\vac\rangle,
 &
 \langle\tPhi_i^a\rvert
 &=\langle\vac\rvert Y_a^i,
 \\
 \lvert\Phi_{ij}^{ab}\rangle
 &=X_{ij}^{ab}\lvert\vac\rangle,
 &
 \langle\tPhi_{ij}^{ab}\rvert
 &=\langle\vac\rvert Y_{ab}^{ij}.
\end{align}
The singles are biorthonormal. The doubles retain the column redundancy
\begin{equation}
 C_{ij}^{ab}=C_{ji}^{ba},
 \qquad
 \tC_{ab}^{ij}=\tC_{ba}^{ji}.
\end{equation}
It is convenient to define the simultaneous column symmetrizer
\begin{equation}
 \mathcal P_{pq}^{rs}A_{pq}^{rs}
 \equiv A_{pq}^{rs}+A_{qp}^{sr}.
\end{equation}
For example,
\begin{equation}
 \mathcal P_{ij}^{ab}A_{ij}^{ab}
 =A_{ij}^{ab}+A_{ji}^{ba}.
\end{equation}
The symmetrizer exchanges both columns; it does not exchange only the upper
or only the lower pair. It is unnormalized and therefore satisfies
\begin{equation}
 \left(\mathcal P_{pq}^{rs}\right)^2
 =2\mathcal P_{pq}^{rs}.
\end{equation}
When it precedes brackets, it acts on every factor and every free index in the
entire bracketed expression. Terms that are already invariant under the
simultaneous exchange are written explicitly rather than placed inside
$\mathcal P$, avoiding an unintended factor of two.
No diagonal normalization factor such as
$(1+\delta_{ab}\delta_{ij})^{-1}$ is introduced.

\subsection{Left and right CI states}

The most general state needed below is
\begin{align}
 \lvert\Psi\rangle
 &=\left(
   C_0+C_i^aX_i^a+\frac12C_{ij}^{ab}X_{ij}^{ab}
   \right)\lvert\vac\rangle,
 \\
 \langle\tPsi\rvert
 &=\langle\vac\rvert\left(
   \tC_0+\tC_a^iY_a^i+\frac12\tC_{ab}^{ij}Y_{ab}^{ij}
   \right).
\end{align}
The left coefficients are kept algebraically independent. When
$\langle\tPsi\rvert=\langle\Psi\rvert$, they are obtained from the right
coefficients through
\begin{align}
 \tC_0&=C_0^{*},
 &
 \tC_a^i&=2(C_i^a)^{*},
 \\
 \tC_{ab}^{ij}
 &=\left(4C_{ij}^{ab}-2C_{ji}^{ab}\right)^{*}.
\end{align}

\subsection{Sigma vectors and reduced density matrices}

The right sigma-vector components are
\begin{align}
 \sigma_0
 &=\langle\vac\rvert H\lvert\Psi\rangle,
 \\
 \sigma_i^a
 &=\langle\tPhi_i^a\rvert H\lvert\Psi\rangle,
 \\
 \sigma_{ij}^{ab}
 &=\langle\tPhi_{ij}^{ab}\rvert H\lvert\Psi\rangle.
\end{align}
The corresponding left action is
\begin{align}
 \tsigma_0
 &=\langle\tPsi\rvert H\lvert\vac\rangle,
 \\
 \tsigma_a^i
 &=\langle\tPsi\rvert H\lvert\Phi_i^a\rangle,
 \\
 \tsigma_{ab}^{ij}
 &=\langle\tPsi\rvert H\lvert\Phi_{ij}^{ab}\rangle.
\end{align}

The spin-summed one- and two-body reduced density matrices are
\begin{align}
 \gamma_p^q
 &\equiv\langle\tPsi\rvert E_{pq}\lvert\Psi\rangle,
 \\
 \Gamma_{pq}^{rs}
 &\equiv\langle\tPsi\rvert e_{pqrs}\lvert\Psi\rangle.
\end{align}
The expressions below are not divided by
$\langle\tPsi\vert\Psi\rangle$. Thus normalization drift remains visible,
as required for applications such as real-time propagation. Only nonzero RDM
blocks are displayed. Two-body blocks related by
$e_{pqrs}=e_{qpsr}$ are represented only once.

\section{Derived algebraic expressions}

\subsection{Configuration-interaction singles (CIS)}

For CIS, $C_{ij}^{ab}=\tC_{ab}^{ij}=0$:
\begin{align}
 \lvert\Psi_{\mathrm{CIS}}\rangle
 &=\left(C_0+C_i^aX_i^a\right)\lvert\vac\rangle,
 \\
 \langle\tPsi_{\mathrm{CIS}}\rvert
 &=\langle\vac\rvert\left(\tC_0+\tC_a^iY_a^i\right).
\end{align}
For a Hermitian bra, $\tC_a^i=2(C_i^a)^*$.

% Closed-shell CIS in the HJO biorthogonal basis
% For a Hermitian bra: \widetilde{C}_{0}=C_{0}^{*}, \widetilde{C}^{i}_{a}=2(C^{a}_{i})^{*}.
% No normalization denominator has been applied.

\subsubsection{Norm}
\begin{align*}
\langle\widetilde{\Psi}\vert\Psi\rangle &= C_{0} \widetilde{C}_{0}   + C^{a}_{i}  \widetilde{C}^{i}_{a}
\end{align*}

\subsubsection{Right sigma vector}
\begin{align*}
\sigma_{0} &= C_{0} E_{\mathrm{ref}}   + 2 C^{a}_{i}  f^{i}_{a} \\
\sigma^{a}_{i} &= C_{0} f^{a}_{i}  \\
  &\quad  + E_{\mathrm{ref}} C^{a}_{i}  \\
  &\quad  + C^{b}_{i}  f^{a}_{b}  \\
  &\quad  + C^{b}_{j}  L^{aj}_{ib}  \\
  &\quad - C^{a}_{j}  f^{j}_{i}
\end{align*}

\subsubsection{Left sigma vector}
\begin{align*}
\widetilde{\sigma}_{0} &= \widetilde{C}_{0} E_{\mathrm{ref}}   + \widetilde{C}^{i}_{a}  f^{a}_{i} \\
\widetilde{\sigma}^{i}_{a} &= E_{\mathrm{ref}} \widetilde{C}^{i}_{a}  \\
  &\quad  + 2 \widetilde{C}_{0} f^{i}_{a}  \\
  &\quad  + \widetilde{C}^{i}_{b}  f^{b}_{a}  \\
  &\quad  + \widetilde{C}^{j}_{b}  L^{ib}_{aj}  \\
  &\quad - \widetilde{C}^{j}_{a}  f^{i}_{j}
\end{align*}

\subsubsection{One-body reduced density matrix}
\begin{align*}
\gamma^{j}_{i} &= 2 C_{0} \widetilde{C}_{0} \delta_{i j}   + 2 \delta_{i j} C^{a}_{k}  \widetilde{C}^{k}_{a}  - C^{a}_{i}  \widetilde{C}^{j}_{a} \\
\gamma^{a}_{i} &= 2 \widetilde{C}_{0} C^{a}_{i} \\
\gamma^{i}_{a} &= C_{0} \widetilde{C}^{i}_{a} \\
\gamma^{b}_{a} &= C^{b}_{i}  \widetilde{C}^{i}_{a}
\end{align*}

\subsubsection{Two-body reduced density matrix}
\begin{align*}
\Gamma^{kl}_{ij} &= -  2 C_{0} \widetilde{C}_{0} \delta_{i l} \delta_{j k}  \\
  &\quad  + 4 C_{0} \widetilde{C}_{0} \delta_{i k} \delta_{j l}  \\
  &\quad -  2 \delta_{i l} \delta_{j k} C^{a}_{i_{0}}  \widetilde{C}^{i_{0}}_{a}  \\
  &\quad  + 4 \delta_{i k} \delta_{j l} C^{a}_{i_{0}}  \widetilde{C}^{i_{0}}_{a} \\
  &\quad + \mathcal{P}^{kl}_{ij}\Bigl[ \delta_{i l} C^{a}_{j}  \widetilde{C}^{k}_{a}  \\
  &\quad -  2 \delta_{i k} C^{a}_{j}  \widetilde{C}^{l}_{a} \Bigr] \\
\Gamma^{ak}_{ij} &= -  2 \widetilde{C}_{0} \delta_{i k} C^{a}_{j}   + 4 \widetilde{C}_{0} \delta_{j k} C^{a}_{i} \\
\Gamma^{jk}_{ai} &= -  C_{0} \delta_{i j} \widetilde{C}^{k}_{a}   + 2 C_{0} \delta_{i k} \widetilde{C}^{j}_{a} \\
\Gamma^{bj}_{ai} &= - C^{b}_{i}  \widetilde{C}^{j}_{a}   + 2 \delta_{i j} C^{b}_{k}  \widetilde{C}^{k}_{a} \\
\Gamma^{bj}_{ia} &= 2 C^{b}_{i}  \widetilde{C}^{j}_{a}  -  \delta_{i j} C^{b}_{k}  \widetilde{C}^{k}_{a}
\end{align*}


\clearpage
\subsection{Configuration-interaction doubles (CID)}

For CID, $C_i^a=\tC_a^i=0$:
\begin{align}
 \lvert\Psi_{\mathrm{CID}}\rangle
 &=\left(C_0+\frac12C_{ij}^{ab}X_{ij}^{ab}\right)
   \lvert\vac\rangle,
 \\
 \langle\tPsi_{\mathrm{CID}}\rvert
 &=\langle\vac\rvert\left(
   \tC_0+\frac12\tC_{ab}^{ij}Y_{ab}^{ij}
   \right).
\end{align}
For a Hermitian bra,
$\tC_{ab}^{ij}=(4C_{ij}^{ab}-2C_{ji}^{ab})^*$.

% Closed-shell CID in the HJO biorthogonal basis
% Column symmetry: C^{ab}_{ij}=C^{ba}_{ji}, \widetilde{C}^{ij}_{ab}=\widetilde{C}^{ji}_{ba}.
% For a Hermitian bra: \widetilde{C}_{0}=C_{0}^{*}, \widetilde{C}^{ij}_{ab}=\left(4C^{ab}_{ij}-2C^{ab}_{ji}\right)^{*}.
% No normalization denominator has been applied.

\subsubsection{Norm}
\begin{align*}
\langle\widetilde{\Psi}\vert\Psi\rangle &= C_{0} \widetilde{C}_{0}   + \frac{1}{2} C^{ab}_{ij}  \widetilde{C}^{ij}_{ab}
\end{align*}

\subsubsection{Right sigma vector}
\begin{align*}
\sigma_{0} &= C_{0} E_{\mathrm{ref}}   + C^{ab}_{ij}  L^{ij}_{ab} \\
\sigma^{ab}_{ij} &= C^{ab}_{kl}  u^{kl}_{ij}  \\
  &\quad  + C_{0} u^{ab}_{ij}  \\
  &\quad  + E_{\mathrm{ref}} C^{ab}_{ij}  \\
  &\quad  + C^{cd}_{ij}  u^{ab}_{cd} \\
  &\quad + \mathcal{P}^{ab}_{ij}\Bigl[ f^{b}_{c}  C^{ac}_{ij}  \\
  &\quad - f^{k}_{j}  C^{ab}_{ik}  \\
  &\quad  + C^{ac}_{ik}  L^{bk}_{jc}  \\
  &\quad - C^{ac}_{ki}  u^{bk}_{jc}  \\
  &\quad - C^{ac}_{kj}  u^{bk}_{ci} \Bigr]
\end{align*}

\subsubsection{Left sigma vector}
\begin{align*}
\widetilde{\sigma}_{0} &= \widetilde{C}_{0} E_{\mathrm{ref}}   + \frac{1}{2} \widetilde{C}^{ij}_{ab}  u^{ab}_{ij} \\
\widetilde{\sigma}^{ij}_{ab} &= E_{\mathrm{ref}} \widetilde{C}^{ij}_{ab}  \\
  &\quad  + 2 \widetilde{C}_{0} L^{ij}_{ab}  \\
  &\quad  + \widetilde{C}^{kl}_{ab}  u^{ij}_{kl}  \\
  &\quad  + \widetilde{C}^{ij}_{cd}  u^{cd}_{ab} \\
  &\quad + \mathcal{P}^{ij}_{ab}\Bigl[ f^{c}_{b}  \widetilde{C}^{ij}_{ac}  \\
  &\quad - f^{j}_{k}  \widetilde{C}^{ik}_{ab}  \\
  &\quad  + \widetilde{C}^{ik}_{ac}  L^{cj}_{kb}  \\
  &\quad - \widetilde{C}^{ik}_{bc}  u^{cj}_{ka}  \\
  &\quad - \widetilde{C}^{ik}_{cb}  u^{cj}_{ak} \Bigr]
\end{align*}

\subsubsection{One-body reduced density matrix}
\begin{align*}
\gamma^{j}_{i} &= 2 C_{0} \widetilde{C}_{0} \delta_{i j}  - C^{ab}_{ik}  \widetilde{C}^{kj}_{ba}   + \delta_{i j} C^{ab}_{kl}  \widetilde{C}^{kl}_{ab} \\
\gamma^{b}_{a} &= C^{bc}_{ij}  \widetilde{C}^{ij}_{ac}
\end{align*}

\subsubsection{Two-body reduced density matrix}
\begin{align*}
\Gamma^{kl}_{ij} &= -  2 C_{0} \widetilde{C}_{0} \delta_{i l} \delta_{j k}  \\
  &\quad  + 4 C_{0} \widetilde{C}_{0} \delta_{i k} \delta_{j l}  \\
  &\quad -  \delta_{i l} \delta_{j k} C^{ab}_{i_{0}i_{1}}  \widetilde{C}^{i_{0}i_{1}}_{ab}  \\
  &\quad  + 2 \delta_{i k} \delta_{j l} C^{ab}_{i_{0}i_{1}}  \widetilde{C}^{i_{0}i_{1}}_{ab}  \\
  &\quad  + C^{ab}_{ij}  \widetilde{C}^{kl}_{ab} \\
  &\quad + \mathcal{P}^{kl}_{ij}\Bigl[ \delta_{i l} C^{ab}_{i_{0}j}  \widetilde{C}^{i_{0}k}_{ab}  \\
  &\quad -  2 \delta_{i k} C^{ab}_{i_{0}j}  \widetilde{C}^{i_{0}l}_{ab} \Bigr] \\
\Gamma^{ab}_{ij} &= -  2 \widetilde{C}_{0} C^{ab}_{ji}   + 4 \widetilde{C}_{0} C^{ab}_{ij} \\
\Gamma^{ij}_{ab} &= C_{0} \widetilde{C}^{ij}_{ab} \\
\Gamma^{bj}_{ai} &= 2 \delta_{i j} C^{bc}_{kl}  \widetilde{C}^{kl}_{ac}  - C^{bc}_{ik}  \widetilde{C}^{kj}_{ca}  - C^{bc}_{ki}  \widetilde{C}^{kj}_{ac} \\
\Gamma^{bj}_{ia} &= -  \delta_{i j} C^{bc}_{kl}  \widetilde{C}^{kl}_{ac}  - C^{bc}_{ki}  \widetilde{C}^{kj}_{ca}   + 2 C^{bc}_{ik}  \widetilde{C}^{kj}_{ca} \\
\Gamma^{cd}_{ab} &= C^{cd}_{ij}  \widetilde{C}^{ij}_{ab}
\end{align*}


\clearpage
\subsection{Configuration-interaction singles and doubles (CISD)}

The CISD states contain both excitation ranks:
\begin{align}
 \lvert\Psi_{\mathrm{CISD}}\rangle
 &=\left(
   C_0+C_i^aX_i^a+\frac12C_{ij}^{ab}X_{ij}^{ab}
   \right)\lvert\vac\rangle,
 \\
 \langle\tPsi_{\mathrm{CISD}}\rvert
 &=\langle\vac\rvert\left(
   \tC_0+\tC_a^iY_a^i+\frac12\tC_{ab}^{ij}Y_{ab}^{ij}
   \right).
\end{align}
For a Hermitian bra, both left-amplitude transformations given in
Section~1.6 are applied.

\begingroup
\setlength{\jot}{1pt}
% Closed-shell CISD in the HJO biorthogonal basis
% Column symmetry: C^{ab}_{ij}=C^{ba}_{ji}, \widetilde{C}^{ij}_{ab}=\widetilde{C}^{ji}_{ba}.
% For a Hermitian bra: \widetilde{C}_{0}=C_{0}^{*}, \widetilde{C}^{i}_{a}=2(C^{a}_{i})^{*}, \widetilde{C}^{ij}_{ab}=\left(4C^{ab}_{ij}-2C^{ab}_{ji}\right)^{*}.
% No normalization denominator has been applied.

\subsubsection{Norm}
\begin{align*}
\langle\widetilde{\Psi}\vert\Psi\rangle &= C_{0} \widetilde{C}_{0}   + \frac{1}{2} C^{ab}_{ij}  \widetilde{C}^{ij}_{ab}   + C^{a}_{i}  \widetilde{C}^{i}_{a}
\end{align*}

\subsubsection{Right sigma vector}
\begin{align*}
\sigma_{0} &= C_{0} E_{\mathrm{ref}}   + C^{ab}_{ij}  L^{ij}_{ab}   + 2 C^{a}_{i}  f^{i}_{a} \\
\sigma^{a}_{i} &= C_{0} f^{a}_{i}  \\
  &\quad  + E_{\mathrm{ref}} C^{a}_{i}  \\
  &\quad  + C^{bc}_{ij}  L^{aj}_{bc}  \\
  &\quad  + C^{b}_{i}  f^{a}_{b}  \\
  &\quad - C^{ab}_{jk}  L^{jk}_{ib}  \\
  &\quad - f^{j}_{b}  C^{ab}_{ji}  \\
  &\quad  + C^{b}_{j}  L^{aj}_{ib}  \\
  &\quad  + 2 f^{j}_{b}  C^{ab}_{ij}  \\
  &\quad - C^{a}_{j}  f^{j}_{i} \\
\sigma^{ab}_{ij} &= C^{ab}_{kl}  u^{kl}_{ij}  \\
  &\quad  + C_{0} u^{ab}_{ij}  \\
  &\quad  + E_{\mathrm{ref}} C^{ab}_{ij}  \\
  &\quad  + C^{cd}_{ij}  u^{ab}_{cd} \\
  &\quad + \mathcal{P}^{ab}_{ij}\Bigl[ C^{a}_{i}  f^{b}_{j}  \\
  &\quad  + C^{c}_{i}  u^{ab}_{cj}  \\
  &\quad  + f^{b}_{c}  C^{ac}_{ij}  \\
  &\quad - C^{a}_{k}  u^{bk}_{ji}  \\
  &\quad - f^{k}_{j}  C^{ab}_{ik}  \\
  &\quad  + C^{ac}_{ik}  L^{bk}_{jc}  \\
  &\quad - C^{ac}_{ki}  u^{bk}_{jc}  \\
  &\quad - C^{ac}_{kj}  u^{bk}_{ci} \Bigr]
\end{align*}

\subsubsection{Left sigma vector}
\begin{align*}
\widetilde{\sigma}_{0} &= \widetilde{C}_{0} E_{\mathrm{ref}}   + \frac{1}{2} \widetilde{C}^{ij}_{ab}  u^{ab}_{ij}   + \widetilde{C}^{i}_{a}  f^{a}_{i} \\
\widetilde{\sigma}^{i}_{a} &= E_{\mathrm{ref}} \widetilde{C}^{i}_{a}  \\
  &\quad  + 2 \widetilde{C}_{0} f^{i}_{a}  \\
  &\quad  + \widetilde{C}^{ij}_{bc}  u^{bc}_{aj}  \\
  &\quad  + \widetilde{C}^{i}_{b}  f^{b}_{a}  \\
  &\quad  + f^{b}_{j}  \widetilde{C}^{ij}_{ab}  \\
  &\quad  + \widetilde{C}^{j}_{b}  L^{ib}_{aj}  \\
  &\quad - \widetilde{C}^{j}_{a}  f^{i}_{j}  \\
  &\quad - \widetilde{C}^{jk}_{ab}  u^{ib}_{jk} \\
\widetilde{\sigma}^{ij}_{ab} &= E_{\mathrm{ref}} \widetilde{C}^{ij}_{ab}  \\
  &\quad  + 2 \widetilde{C}_{0} L^{ij}_{ab}  \\
  &\quad  + \widetilde{C}^{kl}_{ab}  u^{ij}_{kl}  \\
  &\quad  + \widetilde{C}^{ij}_{cd}  u^{cd}_{ab} \\
  &\quad + \mathcal{P}^{ij}_{ab}\Bigl[ - \widetilde{C}^{i}_{b}  f^{j}_{a}  \\
  &\quad  + 2 \widetilde{C}^{i}_{a}  f^{j}_{b}  \\
  &\quad  + f^{c}_{b}  \widetilde{C}^{ij}_{ac}  \\
  &\quad  + \widetilde{C}^{i}_{c}  L^{cj}_{ab}  \\
  &\quad - f^{j}_{k}  \widetilde{C}^{ik}_{ab}  \\
  &\quad - \widetilde{C}^{k}_{a}  L^{ij}_{kb}  \\
  &\quad  + \widetilde{C}^{ik}_{ac}  L^{cj}_{kb}  \\
  &\quad - \widetilde{C}^{ik}_{bc}  u^{cj}_{ka}  \\
  &\quad - \widetilde{C}^{ik}_{cb}  u^{cj}_{ak} \Bigr]
\end{align*}

\subsubsection{One-body reduced density matrix}
\begin{align*}
\gamma^{j}_{i} &= 2 C_{0} \widetilde{C}_{0} \delta_{i j}  \\
  &\quad  + 2 \delta_{i j} C^{a}_{k}  \widetilde{C}^{k}_{a}  \\
  &\quad - C^{ab}_{ik}  \widetilde{C}^{kj}_{ba}  \\
  &\quad  + \delta_{i j} C^{ab}_{kl}  \widetilde{C}^{kl}_{ab}  \\
  &\quad - C^{a}_{i}  \widetilde{C}^{j}_{a} \\
\gamma^{a}_{i} &= 2 \widetilde{C}_{0} C^{a}_{i}  - \widetilde{C}^{j}_{b}  C^{ab}_{ji}   + 2 \widetilde{C}^{j}_{b}  C^{ab}_{ij} \\
\gamma^{i}_{a} &= C_{0} \widetilde{C}^{i}_{a}   + C^{b}_{j}  \widetilde{C}^{ij}_{ab} \\
\gamma^{b}_{a} &= C^{bc}_{ij}  \widetilde{C}^{ij}_{ac}   + C^{b}_{i}  \widetilde{C}^{i}_{a}
\end{align*}

\subsubsection{Two-body reduced density matrix}
\begin{align*}
\Gamma^{kl}_{ij} &= -  2 C_{0} \widetilde{C}_{0} \delta_{i l} \delta_{j k}  \\
  &\quad  + 4 C_{0} \widetilde{C}_{0} \delta_{i k} \delta_{j l}  \\
  &\quad -  \delta_{i l} \delta_{j k} C^{ab}_{i_{0}i_{1}}  \widetilde{C}^{i_{0}i_{1}}_{ab}  \\
  &\quad  + 2 \delta_{i k} \delta_{j l} C^{ab}_{i_{0}i_{1}}  \widetilde{C}^{i_{0}i_{1}}_{ab}  \\
  &\quad -  2 \delta_{i l} \delta_{j k} C^{a}_{i_{0}}  \widetilde{C}^{i_{0}}_{a}  \\
  &\quad  + 4 \delta_{i k} \delta_{j l} C^{a}_{i_{0}}  \widetilde{C}^{i_{0}}_{a}  \\
  &\quad  + C^{ab}_{ij}  \widetilde{C}^{kl}_{ab} \\
  &\quad + \mathcal{P}^{kl}_{ij}\Bigl[ \delta_{i l} C^{ab}_{i_{0}j}  \widetilde{C}^{i_{0}k}_{ab}  \\
  &\quad -  2 \delta_{i k} C^{ab}_{i_{0}j}  \widetilde{C}^{i_{0}l}_{ab}  \\
  &\quad  + \delta_{i l} C^{a}_{j}  \widetilde{C}^{k}_{a}  \\
  &\quad -  2 \delta_{i k} C^{a}_{j}  \widetilde{C}^{l}_{a} \Bigr] \\
\Gamma^{ak}_{ij} &= -  2 \widetilde{C}_{0} \delta_{i k} C^{a}_{j}  \\
  &\quad  + 4 \widetilde{C}_{0} \delta_{j k} C^{a}_{i}  \\
  &\quad  + \delta_{i k} \widetilde{C}^{l}_{b}  C^{ab}_{lj}  \\
  &\quad -  2 \delta_{i k} \widetilde{C}^{l}_{b}  C^{ab}_{jl}  \\
  &\quad -  2 \delta_{j k} \widetilde{C}^{l}_{b}  C^{ab}_{li}  \\
  &\quad  + 4 \delta_{j k} \widetilde{C}^{l}_{b}  C^{ab}_{il}  \\
  &\quad  + \widetilde{C}^{k}_{b}  C^{ab}_{ji}  \\
  &\quad - 2 \widetilde{C}^{k}_{b}  C^{ab}_{ij} \\
\Gamma^{jk}_{ai} &= -  \delta_{i j} C^{b}_{l}  \widetilde{C}^{lk}_{ba}  \\
  &\quad  + 2 \delta_{i k} C^{b}_{l}  \widetilde{C}^{lj}_{ba}  \\
  &\quad -  C_{0} \delta_{i j} \widetilde{C}^{k}_{a}  \\
  &\quad  + 2 C_{0} \delta_{i k} \widetilde{C}^{j}_{a}  \\
  &\quad - C^{b}_{i}  \widetilde{C}^{jk}_{ab} \\
\Gamma^{ab}_{ij} &= -  2 \widetilde{C}_{0} C^{ab}_{ji}   + 4 \widetilde{C}_{0} C^{ab}_{ij} \\
\Gamma^{ij}_{ab} &= C_{0} \widetilde{C}^{ij}_{ab} \\
\Gamma^{bj}_{ai} &= - C^{b}_{i}  \widetilde{C}^{j}_{a}  \\
  &\quad  + 2 \delta_{i j} C^{bc}_{kl}  \widetilde{C}^{kl}_{ac}  \\
  &\quad  + 2 \delta_{i j} C^{b}_{k}  \widetilde{C}^{k}_{a}  \\
  &\quad - C^{bc}_{ik}  \widetilde{C}^{kj}_{ca}  \\
  &\quad - C^{bc}_{ki}  \widetilde{C}^{kj}_{ac} \\
\Gamma^{bj}_{ia} &= 2 C^{b}_{i}  \widetilde{C}^{j}_{a}  \\
  &\quad -  \delta_{i j} C^{bc}_{kl}  \widetilde{C}^{kl}_{ac}  \\
  &\quad -  \delta_{i j} C^{b}_{k}  \widetilde{C}^{k}_{a}  \\
  &\quad - C^{bc}_{ki}  \widetilde{C}^{kj}_{ca}  \\
  &\quad  + 2 C^{bc}_{ik}  \widetilde{C}^{kj}_{ca} \\
\Gamma^{bc}_{ai} &= - \widetilde{C}^{j}_{a}  C^{bc}_{ij}   + 2 \widetilde{C}^{j}_{a}  C^{bc}_{ji} \\
\Gamma^{ci}_{ab} &= C^{c}_{j}  \widetilde{C}^{ij}_{ba} \\
\Gamma^{cd}_{ab} &= C^{cd}_{ij}  \widetilde{C}^{ij}_{ab}
\end{align*}

\endgroup

\clearpage
\appendix
\section{Reproducibility}

The equation fragments included by this document are generated directly from
the validated Drudge scripts. Drudge first produces fully expanded tensors.
A symbolic presentation pass then factors exact Coulomb-exchange pairs through
$L_{rs}^{pq}$ and groups distinct simultaneous column-exchange pairs under
$\mathcal P_{pq}^{rs}$. It re-expands both abbreviations and requires the
difference from the original Drudge tensor to simplify to zero. Neither
abbreviation is introduced by textual manipulation of the LaTeX during normal
generation.

From this directory, regenerate the fragments with
\begin{verbatim}
conda run --no-capture-output -n drudge python generate_latex_fragments.py
\end{verbatim}
and compile the document with
\begin{verbatim}
latexmk -pdf -interaction=nonstopmode -halt-on-error \
  closed_shell_ci_expressions.tex
\end{verbatim}

\end{document}
